\subsubsection{极点留数与多极矩增长：Ramanujan--视界二分律的定量形态}\label{subsubsec:xi-horizon-residue-multipole-temperedness}
附录 \ref{subsubsec:app-horizon-weyl-carath-toeplitz} 已将 RH 等价表述为视界对数导数
\(\mathscr{C}_{\zeta}\) 的 Carath\'eodory 正性，并进一步等价为 Toeplitz--PSD 层级。
同时，离临界线偏移在圆盘半径上的二次径向压缩律由 \eqref{eq:app-offcritical-radius-closed}（亦见定理 \ref{thm:xi-offcritical-quadratic-radial-compression}）给出。
本小节给出三条严格超出上述“等价式 + 位置压缩律”的\emph{定量}结果：离线零点不仅被压入端点薄层，而且在 \(\mathscr{C}_{\zeta}\) 上诱导一条留数强度同阶压缩律；并且该盘内极点在视界多极矩 \(\{c_n\}\) 上强迫指数型增长尾项，从而在矩谱半径意义下给出一条与 Ramanujan temperedness 同型的二分判据。

\paragraph{视界对数导数与多极矩}
沿用 \eqref{eq:app-horizon-carath-def} 的定义：
\[
\mathscr{C}_{\zeta}(w):=-\frac{\Xi_{\zeta}'(s(w))}{\Xi_{\zeta}(s(w))},\qquad
s(w)=\tfrac12+\mathrm{i}\,\mathcal{C}^{-1}(w)\qquad(w\in\mathbb{D}),
\]
并以归一化幂级数
\begin{equation}\label{eq:xi-horizon-carath-series-recall}
\mathscr{C}_{\zeta}(w)=c_0+2\sum_{n\ge 1}c_n w^n
\end{equation}
定义视界多极矩（亦即“可见毛发系数”）。

\begin{theorem}[极点留数公式与“深度—强度”压缩律]\label{thm:xi-horizon-pole-residue-depth-strength}
\leanverified{paper\_xi\_horizon\_pole\_residue\_depth\_strength}
设 \(\Xi_{\zeta}\) 存在零点 \(\rho=\beta+\mathrm{i}\gamma\) 满足 \(0<\beta<\tfrac12\)，并令
\[
\delta:=\tfrac12-\beta>0,\qquad
t_{\rho}:=\gamma+\mathrm{i}\delta\in\mathbb{H},\qquad
w_{\rho}:=\mathcal{C}(t_{\rho})\in\mathbb{D}.
\]
记该零点的重数为 \(m_{\rho}\ge 1\)。
则 \(\mathscr{C}_{\zeta}\) 在 \(w=w_{\rho}\) 处具有简单极点，且其留数满足完全闭式
\begin{equation}\label{eq:xi-horizon-pole-residue-closed}
\boxed{\;
\operatorname*{Res}_{w=w_{\rho}}\mathscr{C}_{\zeta}(w)
=-\;m_{\rho}\,\frac{(1+w_{\rho})^{2}}{2}
=-\;\frac{2m_{\rho}}{(1-\mathrm{i}t_{\rho})^{2}}\,.
\;}
\end{equation}
此外，令端点深度（视界深度）
\(
h(\rho):=1-|w_{\rho}|^{2}
\)
（见 \eqref{eq:app-offcritical-radius-closed}），则留数模长满足“深度—强度”压缩律
\begin{equation}\label{eq:xi-horizon-depth-strength-law}
\boxed{\;
\Bigl|\operatorname*{Res}_{w=w_{\rho}}\mathscr{C}_{\zeta}(w)\Bigr|
=\frac{2m_{\rho}}{\gamma^{2}+(1+\delta)^{2}}
=\frac{m_{\rho}}{2\delta}\,h(\rho)\; }.
\end{equation}
\end{theorem}

\begin{proof}
令 \(F_{\zeta}(w):=\Xi_{\zeta}(s(w))\)（见 \eqref{eq:app-Fzeta-disk}），则 \(F_{\zeta}\) 在 \(w_{\rho}\) 处零点阶数为 \(m_{\rho}\)。
因此对数导数 \(\frac{\dd}{\dd w}\log F_{\zeta}(w)\) 在 \(w_{\rho}\) 处具有留数 \(m_{\rho}\) 的简单极点。
又由链式法则
\[
\frac{\dd}{\dd w}\log F_{\zeta}(w)
=\frac{\Xi_{\zeta}'(s(w))}{\Xi_{\zeta}(s(w))}\,s'(w),
\qquad
\mathscr{C}_{\zeta}(w)=-\frac{1}{s'(w)}\frac{\dd}{\dd w}\log F_{\zeta}(w),
\]
从而
\[
\operatorname*{Res}_{w=w_{\rho}}\mathscr{C}_{\zeta}(w)
=-\frac{m_{\rho}}{s'(w_{\rho})}.
\]
由 \eqref{eq:app-s-of-w} 有 \(s'(w)=\frac{2}{(1+w)^2}\)，故得到第一等式。
另一方面，由 \eqref{eq:app-cayley-disk} 有 \(1+w=\frac{2}{1-\mathrm{i}t}\)，代入得第二等式。
\par
对模长，注意到
\(
|1-\mathrm{i}t_{\rho}|^2=\gamma^2+(1+\delta)^2
\)
从而得到 \(\bigl|\operatorname*{Res}\bigr|=\frac{2m_{\rho}}{\gamma^2+(1+\delta)^2}\)。
再由推论 \ref{cor:app-offcritical-horocycle-busemann} 的恒等式
\(
\delta=\frac{1-|w_{\rho}|^2}{|1+w_{\rho}|^2}
\)
得 \(|1+w_{\rho}|^2=h(\rho)/\delta\)，与
\(
\bigl|\operatorname*{Res}\bigr|=\frac{m_{\rho}}{2}|1+w_{\rho}|^2
\)
合并即得 \eqref{eq:xi-horizon-depth-strength-law}。证毕。
\end{proof}

\begin{corollary}[红移抑制：高高度下留数强度与深度同阶衰减]\label{cor:xi-horizon-residue-redshift}
\leanverified{paper\_xi\_horizon\_residue\_redshift}
在定理 \ref{thm:xi-horizon-pole-residue-depth-strength} 的设定下，若 \(|\gamma|\to\infty\) 且 \(\delta\) 有界，则
\[
h(\rho)\sim \frac{4\delta}{\gamma^2},
\qquad
\Bigl|\operatorname*{Res}_{w=w_{\rho}}\mathscr{C}_{\zeta}(w)\Bigr|\sim \frac{2m_{\rho}}{\gamma^2}.
\]
因此离线零点不仅被压入端点薄层（位置红移），其对视界可观测函数 \(\mathscr{C}_{\zeta}\) 的耦合强度亦在同一 \(\gamma^{-2}\) 标度上被抑制（强度红移）。
\end{corollary}

\begin{theorem}[视界多极矩的 Ramanujan 温和性二分律]\label{thm:xi-horizon-multipole-temperedness-dichotomy}
\leanverified{paper\_xi\_horizon\_multipole\_temperedness\_dichotomy}
令 \(\mathscr{C}_{\zeta}(w)=c_0+2\sum_{n\ge 1}c_n w^n\) 如 \eqref{eq:xi-horizon-carath-series-recall}。
定义矩谱半径（多极矩谱半径）
\[
\mathfrak{r}_{\mathrm{mom}}
:=\limsup_{n\to\infty}|c_n|^{1/n}\in[0,\infty],
\qquad
\Lambda_{\mathrm{mom}}:=\log \mathfrak{r}_{\mathrm{mom}}\in[-\infty,\infty].
\]
则成立如下二分律：
\begin{enumerate}
  \item \textbf{（tempered）}若 RH 成立，则 \(\mathscr{C}_{\zeta}\) 为 Carath\'eodory 函数（推论 \ref{cor:app-horizon-carath-rh}）。
  因而存在有限正测度 \(\mu\) 使 \(c_n=\int_{\TT}\xi^{-n}\,\dd\mu(\xi)\)（见 \eqref{eq:app-horizon-herglotz-disk}--\eqref{eq:app-horizon-carath-series}），从而
  \begin{equation}\label{eq:xi-horizon-cn-bounded}
  \boxed{\;
  |c_n|\le c_0=\mu(\TT)\qquad(\forall n\ge 0)\;}
  \end{equation}
  且 \(\Lambda_{\mathrm{mom}}\le 0\)。
  \item \textbf{（non-tempered）}若 RH 不成立，则 \(\mathscr{C}_{\zeta}\) 在 \(\mathbb{D}\) 内至少有一极点。
  令
  \begin{equation}\label{eq:xi-horizon-rmin-pole}
  r_{\min}:=\min\Bigl\{|a|:\ a\in\mathbb{D},\ \mathscr{C}_{\zeta}\ \text{在}\ a\ \text{处有极点}\Bigr\}\in(0,1).
  \end{equation}
  则 \(\mathscr{C}_{\zeta}\) 在 \(w=0\) 的幂级数收敛半径等于 \(r_{\min}\)，并且
  \begin{equation}\label{eq:xi-horizon-moment-spectral-radius}
  \boxed{\;
  \mathfrak{r}_{\mathrm{mom}}=\frac{1}{r_{\min}}>1,
  \qquad
  \Lambda_{\mathrm{mom}}=-\log r_{\min}>0\; }.
  \end{equation}
  \item \textbf{（闭式指数：由最深盘内极点给出）}若 \(r_{\min}\) 由某个离线零点 \(\rho\) 所诱导的极点 \(w_{\rho}\) 实现，则
  \begin{equation}\label{eq:xi-horizon-Lambda-closed}
  \boxed{\;
  \Lambda_{\mathrm{mom}}
  =-\log|w_{\rho}|
  =\frac12\log\!\left(\frac{\gamma^{2}+(1+\delta)^{2}}{\gamma^{2}+(1-\delta)^{2}}\right)\; }.
  \end{equation}
  且在高高度区间 \(|\gamma|\to\infty\) 且 \(\delta\) 有界时，
  \begin{equation}\label{eq:xi-horizon-Lambda-asymptotic}
  \boxed{\;
  \Lambda_{\mathrm{mom}}\sim \frac{2\delta}{\gamma^{2}}\sim \frac{h(\rho)}{2}\; }.
  \end{equation}
\end{enumerate}
\end{theorem}

\begin{proof}
(1) 由 Carath\'eodory/Herglotz 表示，\(c_n\) 为有限正测度 \(\mu\) 的 Fourier 矩，因此
\(
|c_n|\le \int_{\TT}|\xi^{-n}|\,\dd\mu=\mu(\TT)=c_0
\)，得到 \eqref{eq:xi-horizon-cn-bounded}。
由于 \(\{c_n\}\) 有界，故 \(\limsup |c_n|^{1/n}\le 1\)，从而 \(\Lambda_{\mathrm{mom}}\le 0\)。
\par
(2) RH 不成立时，\(\mathscr{C}_{\zeta}\) 在 \(\mathbb{D}\) 内出现极点（见推论 \ref{cor:app-horizon-carath-rh} 的证明末段与命题 \ref{prop:app-horizon-pole-imprint}）。
由解析函数的一般理论，幂级数在 \(w=0\) 的收敛半径等于到最近奇点的距离；因而收敛半径等于 \(r_{\min}\)。
再由 Cauchy--Hadamard 定理，\(\frac{1}{r_{\min}}=\limsup |2c_n|^{1/n}=\limsup |c_n|^{1/n}\)，得到 \eqref{eq:xi-horizon-moment-spectral-radius}。
\par
(3) 由 \eqref{eq:app-offcritical-radius-closed} 得
\(
|w_{\rho}|^2=\frac{\gamma^2+(1-\delta)^2}{\gamma^2+(1+\delta)^2}
\)，故 \(-\log|w_{\rho}|\) 具有闭式 \eqref{eq:xi-horizon-Lambda-closed}。
又由 \(h(\rho)=1-|w_{\rho}|^2\to 0\)（\(|\gamma|\to\infty\)、\(\delta\) 有界）以及恒等式
\(
-\log|w_{\rho}|=-\tfrac12\log(1-h(\rho))
=\tfrac12 h(\rho)+O(h(\rho)^2)
\)，并用 \(h(\rho)\sim 4\delta/\gamma^2\)（由 \eqref{eq:app-offcritical-radius-closed}）即可得 \eqref{eq:xi-horizon-Lambda-asymptotic}。证毕。
\end{proof}

\begin{corollary}[亚指数增长证书：多极矩的 Ramanujan 归约]\label{cor:xi-horizon-subexp-growth-implies-rh}
若
\[
\limsup_{n\to\infty}\frac{1}{n}\log|c_n|\le 0,
\]
则 RH 成立。
等价地，任意 \( |c_n|\le e^{o(n)} \) 的上界（亚指数增长/不出现正指数增长）均推出 RH。
\end{corollary}

\begin{proof}
若 RH 不成立，则由定理 \ref{thm:xi-horizon-multipole-temperedness-dichotomy} 第 (2) 点有
\(
\limsup_{n\to\infty}\frac{1}{n}\log|c_n|
=\Lambda_{\mathrm{mom}}>0
\)，与假设矛盾；故 RH 成立。证毕。
\end{proof}

\begin{theorem}[单一离线极点对多极矩的强迫几何尾项]\label{thm:xi-horizon-forced-geometric-tail}
在定理 \ref{thm:xi-horizon-pole-residue-depth-strength} 的设定下，记
\(
R_{\rho}:=\operatorname*{Res}_{w=w_{\rho}}\mathscr{C}_{\zeta}(w)
\neq 0
\)。
则存在在 \(\mathbb{D}\) 内解析的函数 \(H_{\rho}\) 使
\begin{equation}\label{eq:xi-horizon-pole-decomposition}
\boxed{\;
\mathscr{C}_{\zeta}(w)=\frac{R_{\rho}}{w-w_{\rho}}+H_{\rho}(w)\qquad(w\in\mathbb{D}).\;}
\end{equation}
将 \(H_{\rho}\) 在 \(w=0\) 处展开为 \(H_{\rho}(w)=h_0+2\sum_{n\ge 1}h_n w^n\)，则对一切 \(n\ge 1\) 有系数恒等式
\begin{equation}\label{eq:xi-horizon-cn-forced-tail}
\boxed{\;
c_n=-\frac{R_{\rho}}{2\,w_{\rho}^{n+1}}+h_n\qquad(n\ge 1).\;}
\end{equation}
特别地，
\[
\limsup_{n\to\infty}|c_n|^{1/n}\ \ge\ \frac{1}{|w_{\rho}|}.
\]
\leanverified{paper\_xi\_horizon\_forced\_geometric\_tail}
\end{theorem}

\begin{proof}
由 \(\mathscr{C}_{\zeta}\) 在 \(\mathbb{D}\setminus\{w_{\rho}\}\) 解析且在 \(w_{\rho}\) 处仅有一个主部，可定义
\(
H_{\rho}(w):=\mathscr{C}_{\zeta}(w)-\frac{R_{\rho}}{w-w_{\rho}}
\)，则 \(H_{\rho}\) 在 \(\mathbb{D}\) 内解析，从而得到 \eqref{eq:xi-horizon-pole-decomposition}。
对 \(\frac{R_{\rho}}{w-w_{\rho}}=-\frac{R_{\rho}}{w_{\rho}}\cdot\frac{1}{1-w/w_{\rho}}
=-\sum_{n\ge 0}R_{\rho}\,w_{\rho}^{-(n+1)}w^n\) 的几何级数展开与 \eqref{eq:xi-horizon-carath-series-recall} 对照即可得到 \eqref{eq:xi-horizon-cn-forced-tail}；末式由 Cauchy--Hadamard 直接推出。证毕。
\end{proof}

\begin{corollary}[强度红移 \(\times\) 多极阶放大：一条显式阈值律]\label{cor:xi-horizon-redshift-growth-threshold}
在定理 \ref{thm:xi-horizon-forced-geometric-tail} 的设定下，记
\(\Lambda(\rho):=-\log|w_{\rho}|>0\)。
则极点所强迫的几何尾项幅度为
\begin{equation}\label{eq:xi-horizon-forced-term-amplitude}
\boxed{\;
\frac{|R_{\rho}|}{2}\,|w_{\rho}|^{-(n+1)}
=\frac{m_{\rho}}{\gamma^{2}+(1+\delta)^{2}}\,
\exp\!\bigl((n+1)\Lambda(\rho)\bigr)\; }.
\end{equation}
在高高度区间 \(|\gamma|\to\infty\) 且 \(\delta\) 有界时，有
\[
\frac{|R_{\rho}|}{2}\,|w_{\rho}|^{-(n+1)}
\asymp
\gamma^{-2}\exp\!\bigl(n\,\Lambda(\rho)\bigr),
\qquad
\Lambda(\rho)\asymp \delta/\gamma^{2},
\]
从而离线零点在视界多极矩上同时表现为：强度层面 \(\asymp\gamma^{-2}\) 的红移抑制，与多极阶 \(n\) 方向的指数阈值放大（其指数同样被端点薄层压缩）。
\end{corollary}

\begin{remark}[非温和性参数的“视界语义”]\label{rem:xi-horizon-nontempered-parameter-semantics}
定理 \ref{thm:xi-horizon-multipole-temperedness-dichotomy} 把 RH 的否证机制压缩为一个可计算的“矩谱半径”参数：
\(\mathfrak{r}_{\mathrm{mom}}>1\) 等价于盘内极点的进入，亦等价于多极矩不可避免的指数尾项；
与之相对，RH 下 \(\mathfrak{r}_{\mathrm{mom}}\le 1\) 且多极矩一致有界，从而在本文语义中对应“tempered/unitary 的视界谱对象”。
在此意义上，\(\Lambda_{\mathrm{mom}}\) 可被视作一条规范的非温和性指数，其闭式压缩 \(\Lambda_{\mathrm{mom}}\sim h(\rho)/2\)（高高度）将“非温和性”与端点薄层几何直接对齐。
\end{remark}

\endinput
